\documentclass[11pt,a4paper]{article}
\usepackage[utf8]{inputenc}
\usepackage[T1]{fontenc}
\usepackage{amsmath, amssymb, amsthm}
\usepackage{geometry}
\usepackage{hyperref}
\usepackage{graphicx}

\geometry{margin=1in}

\title{Formal Security Proofs for Schnorr-based zk-KYC in Permissioned Blockchains}
\author{Real World Assets (RWA) Tokenization Project}
\date{\today}

\newtheorem{theorem}{Theorem}
\newtheorem{lemma}{Lemma}
\newtheorem{definition}{Definition}

\begin{document}

\maketitle

\begin{abstract}
This document provides the formal mathematical proofs for the Non-Interactive Zero-Knowledge (NIZK) identity protocol utilized in the RWA Tokenization platform. We demonstrate that the protocol, based on the Schnorr signature scheme with the Fiat-Shamir heuristic over the secp256k1 elliptic curve, satisfies Completeness, Soundness, and Zero-Knowledge. Furthermore, we formalize its alignment with the Markets in Crypto-Assets (MiCA) regulation, specifically Article 68.
\end{abstract}

\section{Introduction}
To comply with regulatory standards such as KYC/AML and MiCA without compromising user privacy, our system employs a zk-KYC protocol. A central authority (the platform) issues an eligibility credential (a signed claim over the user's public key $Y$). To perform a regulated action, the user provides a NIZK proof of knowledge of the private key $x$ corresponding to $Y$, bound to a specific context (purpose and nullifier).

Let $\mathbb{G}$ be the secp256k1 elliptic curve group of prime order $N$ with generator $G$.
The user's key pair is $(x, Y)$ where $x \in \mathbb{Z}_N^*$ and $Y = xG$.

\section{The Protocol ($\Pi_{zk-KYC}$)}
\subsection{Setup and Prove}
Given context $C \in \{0, 1\}^*$:
\begin{enumerate}
    \item Prover samples $r \xleftarrow{\$} \mathbb{Z}_N^*$.
    \item Prover computes the commitment $R = rG$.
    \item Prover computes the challenge $c = H(R_x \parallel R_y \parallel Y_x \parallel Y_y \parallel C) \pmod N$, where $H$ is SHA-256.
    \item Prover computes the response $s = (r + c \cdot x) \pmod N$.
    \item The proof is $\pi = (R_x, R_y, s)$.
\end{enumerate}

\subsection{Verify}
Given proof $\pi = (R_x, R_y, s)$, public key $Y$, and context $C$:
\begin{enumerate}
    \item Reconstruct $R \in \mathbb{G}$ from $(R_x, R_y)$. If invalid, output \texttt{false}.
    \item Compute $c = H(R_x \parallel R_y \parallel Y_x \parallel Y_y \parallel C) \pmod N$.
    \item Verify that $sG = R + cY$. If equal, output \texttt{true}, else \texttt{false}.
\end{enumerate}

\section{Formal Security Proofs}

\subsection{Completeness}
\begin{theorem}[Completeness]
An honest prover with knowledge of private key $x$ such that $Y = xG$ will always produce a proof $\pi$ that is accepted by the verifier.
\end{theorem}
\begin{proof}
By definition of the honest prover, $s = r + cx \pmod N$.
The verifier checks if $sG = R + cY$.
Substituting $s$:
\[ sG = (r + cx)G = rG + cxG \]
Since $R = rG$ and $Y = xG$, we have:
\[ sG = R + cY \]
Thus, the verification equation always holds.
\end{proof}

\subsection{Soundness}
\begin{theorem}[Soundness]
If a Probabilistic Polynomial-Time (PPT) adversary $\mathcal{A}$ can forge a valid proof $\pi$ for a public key $Y$ without knowledge of $x$ with non-negligible probability $\epsilon$, then $\mathcal{A}$ can be used to solve the Discrete Logarithm Problem (DLP) in $\mathbb{G}$.
\end{theorem}
\begin{proof}
We rely on the Forking Lemma. Suppose $\mathcal{A}$ is an algorithm that outputs a valid proof $(R, c, s)$ for $Y$. We construct an extractor $\mathcal{E}$ that interacts with $\mathcal{A}$ and controls the random oracle $H$.

\begin{enumerate}
    \item $\mathcal{E}$ runs $\mathcal{A}$ with random tape $\rho$. $\mathcal{A}$ queries the random oracle for $H(R, Y, C)$ and eventually outputs a valid proof $(R, c_1, s_1)$ where $s_1G = R + c_1Y$.

    \item $\mathcal{E}$ rewinds $\mathcal{A}$ to the exact state before the query $H(R, Y, C)$, keeping the same random tape $\rho$. $\mathcal{E}$ replies to the query with a new random challenge $c_2 \neq c_1$.

    \item By the Forking Lemma, with probability $\geq \epsilon^2 / q_H$ (where $q_H$ is the number of hash queries), $\mathcal{A}$ outputs a second valid proof $(R, c_2, s_2)$ such that $s_2G = R + c_2Y$.

    \item $\mathcal{E}$ now has two valid transcripts:
\end{enumerate}
\[ s_1G = R + c_1Y \]
\[ s_2G = R + c_2Y \]

Subtracting the two equations yields:
\[ (s_1 - s_2)G = (c_1 - c_2)Y \]
\[ (s_1 - s_2)G = (c_1 - c_2)xG \]

Since $c_1 \neq c_2$, $(c_1 - c_2)$ is invertible modulo $N$. $\mathcal{E}$ extracts $x$:
\[ x = (s_1 - s_2)(c_1 - c_2)^{-1} \pmod N \]

Thus, $\mathcal{E}$ has solved the DLP on secp256k1. Assuming the DLP is hard in $\mathbb{G}$, no such PPT adversary $\mathcal{A}$ can exist. Therefore, the protocol is computationally sound.
\end{proof}

\subsection{Zero-Knowledge}
\begin{theorem}[Zero-Knowledge]
The protocol $\Pi_{zk-KYC}$ leaks no information about the private key $x$, as transcripts can be perfectly simulated without knowledge of $x$.
\end{theorem}
\begin{proof}
We construct a PPT simulator $\mathcal{S}$ that, given only the public key $Y$ and context $C$, produces a valid transcript indistinguishable from one produced by an honest prover interacting with the random oracle.

\textbf{Simulator $\mathcal{S}(Y, C)$}:
\begin{enumerate}
    \item Sample $s \xleftarrow{\$} \mathbb{Z}_N^*$.
    \item Sample $c \xleftarrow{\$} \mathbb{Z}_N^*$.
    \item Compute $R = sG - cY$.
    \item Program the random oracle such that $H(R, Y, C) = c$. (If this causes a collision in the oracle, abort. Since $R$ is uniformly distributed, the probability of collision is $q_H / N$, which is negligible).
    \item Output the transcript $\pi = (R, s)$.
\end{enumerate}

In the real protocol, $r$ is chosen uniformly, so $R = rG$ is uniform in $\mathbb{G}$, and $c = H(\dots)$ is uniform in $\mathbb{Z}_N^*$. Thus, $s = r + cx$ is uniform in $\mathbb{Z}_N^*$.
In the simulated protocol, $s$ and $c$ are chosen uniformly, and $R$ is uniquely determined by $R = sG - cY$, making it uniformly distributed in $\mathbb{G}$.
The joint distribution $(R, c, s)$ produced by $\mathcal{S}$ is statistically identical to the distribution produced by the real prover. Thus, the proof leaks zero knowledge about $x$.
\end{proof}

\section{Alignment with MiCA Regulation}
Article 68 of the Markets in Crypto-Assets (MiCA) regulation mandates strict KYC (Know Your Customer) and AML (Anti-Money Laundering) checks for asset transfers, while ensuring data minimization under GDPR.

\begin{itemize}
    \item \textbf{Verification without Disclosure:} The credential acts as a claim $\mathcal{C}_{issuer}(Y, \{kyc\_ok, \dots\})$. The ZKP proves possession of $Y$ without revealing $x$ or identifying metadata on the ledger.
    \item \textbf{Non-Repudiation \& Replay Prevention:} The context $C = \text{purpose} \parallel \text{nullifier}$ guarantees the proof cannot be replayed for another transaction.
    \item \textbf{Auditability:} Regulators can verify the mathematical correctness of the state transitions on the ledger independently using $\pi$, while relying on the off-chain mapping of $Y \leftrightarrow \text{Real Identity}$ maintained by the regulated issuer (e.g., BNP Paribas) when investigation is required.
\end{itemize}

\end{document}
